% Auxiliary-study integration draft; strength evaluation is pending.
% Not included by manuscript.tex yet.
\subsection{Auxiliary prediction and actor representations}
Auxiliary tasks supply additional learning signals to a policy's representation.
UNREAL combines reinforcement learning with auxiliary control and prediction
objectives \cite{jaderberg2016auxiliary}. Phasic Policy Gradient alternates policy
optimization with auxiliary representation learning and penalizes changes to the
policy during the auxiliary phase \cite{cobbe2021phasic}. Our auxiliary study
uses this separation to compare terminal-outcome and terminal-margin prediction
while retaining match-win PPO and the deployed action heads.

Distributional reinforcement learning predicts a return distribution through
Bellman updates \cite{bellemare2017distributional}. Our target is instead a
completed match's clipped score margin, learned by supervised likelihood.
Imani and White connect benefits of distributional regression losses to their
optimization properties \cite{imani2018distributional}; Lyle et al. examine the
roles of nonlinear approximation, early-layer learning and categorical outputs
\cite{lyle2019comparative}. We keep the categorical head and its initialization
identical across outcome and margin supervision. The comparison changes target
resolution within the loss/draw/win groups, including its resulting gradient
effects, and measures the impact on match play.

Garcin et al. show that separated actor and critic representations can specialize
in different information, and that auxiliary objectives and data diversity
affect their performance \cite{garcin2025interplay}. This motivates evaluating
what a terminal-score prediction task contributes to action selection. Our
fixed-state diagnostics measure action disagreement and entropy; complete
paired matches measure playing strength. These two measurements distinguish
changes in the policy from improvements in its decisions.
